\begin{theorem}[统一实现假设下 Lee--Yang digit 集的刚性四元性]\label{thm:killo-leyang-unified-digit-rigidity}
\leanverified{paper\_killo\_leyang\_unified\_digit\_rigidity}
假设猜想 \ref{conj:killo-godel-leyang-unified-realization} 成立。
进一步假设该 \(2^L\)-进仿射实现对 Lee--Yang 分支塔是忠实且无重叠的，具体地说：记
\[
\Phi_d(x):=2^Lx+d\qquad (d\in D),
\]
并对每个长度 \(n\) 的 digit 词 \(w=(d_1,\dots,d_n)\in D^n\) 定义复合
\[
\Phi_w:=\Phi_{d_n}\circ\cdots\circ\Phi_{d_1}.
\]
若存在五个初始种子 \(z_1,\dots,z_5\in T_2(E_{\min})\)，使得对每个 \(n\ge 0\) 都有不交分解
\[
R(y_n)=\bigsqcup_{i=1}^{5}\{\Phi_w(z_i):\ w\in D^n\},
\]
并且对每个固定的 \(i\)，映射
\[
D^n\to R(y_n),\qquad w\mapsto \Phi_w(z_i)
\]
为单射，则必有
\[
\card{D}=4.
\]
从而相应的符号地址系统每步恰引入两比特新信息；其拓扑熵为
\[
h_{\mathrm{top}}=\log 4=2\log 2.
\]
进一步，若记与该无重叠地址系统对应的极限地址集为 \(K_{\mathrm{LY}}\subset T_2(E_{\min})\)，则
\[
\dim_H(K_{\mathrm{LY}})=\dim_M(K_{\mathrm{LY}})=\frac{\log 4}{L\log 2}=\frac{2}{L}.
\]
\end{theorem}
\begin{proof}
由假设，对每个 \(n\ge 0\)，第 \(n\) 层分支点总数既等于
\[
\card{R(y_n)}=5\,\card{D}^n,
\]
又由命题 \ref{prop:killo-leyang-branchset-hypercube-families} 满足
\[
\card{R(y_n)}=5\cdot 4^n.
\]
比较两式并对所有 \(n\) 成立，只能得到
\[
\card{D}^n=4^n\qquad (\forall n\ge 0),
\]
故
\[
\card{D}=4.
\]

于是长度 \(n\) 的地址词数恰为 \(4^n\)，故相应的符号地址系统是四符号满移位，其拓扑熵为
\[
h_{\mathrm{top}}=\log\card{D}=\log 4=2\log 2.
\]

再看极限地址集 \(K_{\mathrm{LY}}\)。
由无重叠假设，长度 \(n\) 的 cylinder 两两不交，且每条 cylinder 的尺度为 \(2^{-Ln}\)。
因此 \(K_{\mathrm{LY}}\) 在第 \(n\) 层由 \(4^n\) 个半径 \(2^{-Ln}\) 的不交球组成。
经典 Moran 维数公式遂给出
\[
\dim_H(K_{\mathrm{LY}})=\dim_M(K_{\mathrm{LY}})
=
\lim_{n\to\infty}\frac{\log 4^n}{Ln\log 2}
=
\frac{\log 4}{L\log 2}
=
\frac{2}{L}.
\]
证毕。
\end{proof}

\endinput
